\scanpage{316}
\mspage{15}
\uline{Remarques} \unsure{a)|o)} On trouve \unsure{essentiellt|essentiellement}
les $=$ objets ($L_*$ ou $L_\cdot$) pour exprimer
des \uline{opérations} sur invariants additifs
contravariants, ou sur invariants \inserted{``additifs''} covariants
(du type envisagé) --- c'était
clair à priori à cause du passage
de $A$ à $B$ par \inserted{passage à la} catégorie opposée
--- mais l'op. définie par $L_*$ dépend
de $L_*$ de façon contravariante dans le
premier cas, covariante dans le deuxième.

b) On a en principe résolu, et de façon
particulièrement large, la question initiale
de \deleted{\unsure{ce qu'on peut faire sur}}\inserted{déterminer les constructions} possibles sur
invariants contravariants ou covariants
``additifs'' du type envisagé. On trouve \deleted{\unsure{une}}
\inserted{$\unsure{=}$} \uline{grosse} catégorie, $\mathrm{Ab}_*$ ou $\mathrm{Ab}_\cdot$ (\deleted{\unsure{que}}
\unsure{toujours quitte} de celle-ci à l'opposée). Cela
tient au fait que l'on a admis
des catégories \inserted{$A$} de valeurs \inserted{$B$ ou $A$} assez spéciales,
\unsure{encore} où on peut effectuer \uline{ttes}
les $\varprojlim$, ou ttes les $\varinjlim$ --- ce
qui a pour effet que ces $=$ opérations
peuvent s'effectuer sur les \deleted{\unsure{invariants}}\deleted{\unsure{opérations}}
invariants, et op sur les invariants.

On trouve que lorsqu'on est plus
restrictif sur les opérations qu'on
se permet, \deleted{\unsure{(p.ex. si}}\inserted{donc} moins restrictif
sur $A$ et $B$, on trouve des catégories
correspondantes d'opérations (ou constructions)
sur invariants plus petites --- \deleted{\unsure{alors}}
en $=$ temps qu'on trouve des cat.
``universelles'' plus petites \deleted{\unsure{à valeurs}}\inserted{types}
invariants additifs contra- ou cov.
possibles. \unsure{\&} De façon précise, il doit
résulter des principes généraux (que je
n'ai pas détaillés) que si \deleted{$\unsure{\mathbb C_*}$} $F_o:$ \deleted{$\unsure{\widehat{\boldsymbol\Delta}}$}\inserted{$\boldsymbol\Delta^o\longrightarrow B_o$}
est l'objet de
$\underbrace{((\mathrm{Ab}_*)^o)}_{\deleted{\unsure{B}}\,\inserted{A^o}}{}_*
\simeq({\mathrm{Ab}^{\deleted{\unsure{*}}}_*}^{*}_{\deleted{\unsure{*}}})^o$
qui exprime l'objet $\frac12$ simplicial
additif universel, alors la \deleted{\unsure{sous}} catégorie
cherchée d'opérations [correspondant à
un paquet $\unsure{\mathcal L|L}$ de cat. d'indices par
rapport auxquelles on se propose de
prendre soit des $\varprojlim$ --- dans des cat.
$B$ où on regarde des inv. de nature
contravariante --- soit des $\varinjlim$ --- dans
le cas de cat. $A$ où on regarde les inv. de
nature cov. ] est la sous-catégorie \uline{pleine}
de $B_o=(\mathrm{Ab}_*)^o$ engendrée par $F_o(\Delta_*^n)$ grâce
aux $\unsure{\underline{\lim}|\varprojlim}$ de type $\unsure{\mathcal L|L}$, et \deleted{\unsure{comme}}\inserted{correspond}
donc à la ss-catégorie pleine de $\mathrm{Ab}_*$
engendrée par \inserted{les} $G_o(\Delta^n)\deleted{\unsure{\vdash}}$ \inserted{grâce aux} $\varprojlim$ de type $\unsure{\mathcal L|L}$.
\margin{$\unsure{\mathcal U|U}_{\unsure{\mathcal L|L}}$}
\margin{317}
